% ========================================================================================
% CAPITOLO 6 - INTERFACCIA SERIALE
% ========================================================================================
\chapter{Interfaccia seriale}

% ----------------------------------------------------------------------------------------
\section{Esercizio 10: Comunicazione seriale fra due unità}
% ----------------------------------------------------------------------------------------

\subsection{Esercizio 10}

\begin{traccia}
	Partendo dall'implementazione fornita dalla Digilent di un dispositivo UART-RS232
	(componente RS232RefComp.vhd), progettare, implementare e simulare in VHDL un sistema
	composto da 2 unità A e B che condividono lo stesso segnale di clock e comunicano tra loro
	mediante interfaccia seriale. Il sistema A contiene una ROM di 8 locazioni da 1 byte
	ciascuno, un contatore CONT\_A per scandire le locazioni della ROM e una UART\_A, mentre il
	sistema B contiene una memoria MEM di 8 locazioni da 1 byte ciascuno, un contatore CONT\_B
	per scandire le locazioni della MEM e una UART\_B. Quando un segnale WR viene asserito
	nell'unità A, viene prelevato un byte dalla ROM e inviato all'unità B, che dovrà riceverlo
	e salvarlo in MEM.
\end{traccia}

\progetto
\todo{Interfaccia di \ent{RS232RefComp} (tabella \texttt{porte}), schema a blocchi di A e B, automi.}

\implementazione
\todo{Codice delle unità A e B.}

\simulazione
\todo{Frame seriale completo (start, dati, stop) e verifica di MEM.}

\subsection{Esercizio 10bis (solo 9 CFU)}

\begin{traccia}
	Dopo aver simulato il comportamento del sistema, lo si implementi su board usando un
	bottone per il segnale di WR in A, un bottone per il segnale di RD in B, e i display per la
	visualizzazione del dato correntemente trasmesso e memorizzato in MEM. Si testi l'errore di overrun.
\end{traccia}

\sintesiboard
\todo{Top-level, vincoli, dimostrazione dell'overrun (due WR senza RD).}
